% NeurIPS 2026 Concept & Feasibility submission.
%
% Build:
%   cd paper && latexmk -pdf paper.tex
%
% At submission time, replace the placeholder \documentclass with the actual
% NeurIPS style file (download `neurips_2026.sty` and friends from the
% conference site and `\usepackage{neurips_2026}`). The body below is written
% to be drop-in compatible: section structure, figure / table layout, and
% citation style match the NeurIPS template's conventions.
\documentclass{article}

\usepackage[utf8]{inputenc}
\usepackage[T1]{fontenc}
\usepackage[disable]{microtype}
\usepackage[margin=1in]{geometry}
\usepackage{graphicx}
\usepackage{booktabs}
\usepackage{amsmath, amssymb, amsthm}
\usepackage{xcolor}
\usepackage[hidelinks]{hyperref}
\usepackage{url}

% No cleveref dependency -- use \cref-like macros that just delegate to \ref.
\providecommand{\Cref}[1]{Section~\ref{#1}}
\providecommand{\cref}[1]{section~\ref{#1}}

% Theorem styles (informal lemma boxes for the C&F track feel).
\newtheorem{lemma}{Lemma}

% Lightweight code formatting; NeurIPS style ships its own \texttt-style helper.
\newcommand{\code}[1]{\texttt{#1}}

\title{Where the complex inductive bias lives:\\
an architecture-level prior and a separable optimization-level\\
robustness asymmetry, characterized at small scale}

\author{%
  Ashutosh Kumar\\
  \small{Neural Networks in Complex Spaces project}\\
  \small{\url{github.com/ashu1069/Neural-Networks-in-Complex-Spaces}}%
}

\date{}

\begin{document}
\maketitle

\begin{abstract}
The intuition that complex-valued neural networks (CVNNs) help on
phase-bearing signals is widely held and rarely measured carefully. We
demonstrate, at the scale where a mechanism analysis remains tractable,
that the complex inductive bias on IQ classification \emph{decomposes}
into two separable effects with distinct origins. (i) An
\textbf{architectural} prior: \code{ComplexConv1d} is exactly
phase-equivariant (proof in two lines); the magnitude-logit head makes
the network phase-invariant at the readout. (ii) An
\textbf{optimization-level} asymmetry: under matched-shared-trial
hyperparameter selection, real-valued baselines at identical capacity
exhibit a step-1 explosion-into-dead-region failure mode on the
classifier head that complex networks do not. Per-step gradient
telemetry across $5 \times 4 \times 3 = 60$ instrumented runs locates
the explosion specifically on \code{head.weight} (50--80$\times$ larger
gradient at step 1 for real baselines than for complex) and shows it
produces $3/9$ dead seeds (final loss = $\ln 3$, the 3-class chance
level) under three of the five activations and zero dead seeds under
the other two. The same threshold behavior replicates on synthetic
AWGN-only data with the same matched-trial hyperparameters --- the
mechanism is hp-driven, not data-driven. We argue that conflating these
two effects produces the inflated apparent CVNN gaps in some prior
work, and that the appropriate framing for CVNN evaluation is to
control for both effects independently. The 1-D scope, three-class
subset of RadioML 2018.01A, and modest hardware are deliberate: the
scope is chosen to support a complete mechanism analysis with
end-to-end reproducibility (manifests, dirty-bit tracking, structural
CI guards on artifacts). Code, sweep harness, telemetry traces, and
figure scripts are released; full-archive scale-up is parameterized via
a single CLI flag and is signposted future work.
\end{abstract}

\section{Concept and feasibility scope}
\label{sec:scope}

\subsection{What this paper claims and does not claim}

We make three concrete claims, each at a defensible scale and with
explicit reproducibility:

\begin{enumerate}
  \item \textbf{Architectural claim (formal).}
    \code{ComplexConv1d} is phase-equivariant in the bias-free case
    (\Cref{lem:phase-equivariance}); the magnitude-logit head is
    phase-invariant. A real \code{Conv1d} on stacked $(I, Q)$ channels
    does not get this for free --- it must learn rotation invariance
    from data via a $2\times 2$ channel mix the complex network gets by
    construction.
  \item \textbf{Empirical claim.} Complex-valued networks beat
    capacity-matched real-valued ones on synthetic IQ modulation
    classification (\Cref{sec:synthetic-rf}) and on a 3-class subset of
    RadioML 2018.01A~\cite{oshea2018radioml}
    (\Cref{sec:radioml}) under matched-budget evaluation. Conversely,
    on a task with no temporal phase structure (1-D phase
    classification, \Cref{sec:phase-classification}), the complex
    inductive bias offers no measurable advantage. Both directions of
    evidence are reported.
  \item \textbf{Mechanism claim.} The apparent magnitude of the
    RadioML win is activation-dependent
    (\Cref{sec:radioml-ablation}) because matched-shared-trial
    selection picks high-lr regimes under some activations that real
    baselines cannot tolerate. Per-step gradient telemetry
    (\Cref{sec:radioml-mechanism}) localizes the failure to a step-1
    head-gradient explosion that puts a fraction of real-baseline seeds
    into a dead-ReLU region they cannot escape; the threshold behavior
    replicates on synthetic data. Phase equivariance of the
    \emph{layer} is necessary for the architectural prior, but
    activation-level phase equivariance does not predict the empirical
    robustness asymmetry --- locating the asymmetry one level lower, in
    $(\mathrm{Re}, \mathrm{Im})$ parameter coupling at the optimization
    level, is what makes the picture self-consistent
    (\Cref{sec:phase-equivariance}).
\end{enumerate}

\noindent We do \emph{not} claim:

\begin{itemize}
  \item That CVNNs beat real-valued networks at full RadioML 2018.01A
    scale (24 modulations $\times$ 26 SNR levels $\times$ 1024-sample
    length). The infrastructure for that experiment is in place
    (\code{experiments/rf/sweep\_radioml.py --preset full}), but the
    run is deferred. The mechanism analysis is the contribution;
    scale-up is signposted future work.
  \item That the inductive bias survives in 2-D regimes (audio STFT,
    MRI). \code{ComplexConv2d} is not implemented; the project is 1-D
    only.
  \item That CReLU activations generalize across signal-processing
    settings. The activation ablation
    (\Cref{sec:radioml-ablation}) is part of the contribution
    precisely because activation choice changes the appearance (though
    not the direction) of CVNN-vs-real gaps.
\end{itemize}

\subsection{Why feasibility-scale is the right scope}

A mechanism-level study requires the ability to (a) instrument every
parameter's gradient at every step, (b) run hundreds of training
configurations within a working budget, and (c) confirm cross-task
generalization with a second data distribution. All three are tractable
at the scope chosen --- 60 instrumented runs per data source, a
re-runnable budgeted sweep, a synthetic stand-in matching the real-data
benchmark on architecture and hyperparameter regime. They are not
tractable for paper deadlines at 24-class $\times$ 1024-sample
$\times$ 100-seed scale. We treat the smaller scope as a positive
feature: it lets us report what is happening, not just that something
is happening.

\subsection{Reproducibility as a first-class concern}

Every empirical claim above is backed by a committed result manifest
that records the Python / PyTorch / OS / device / dtype, every seed
used, the exact \code{git\_commit} of the producing code, and a
\code{git\_dirty} boolean. CI surfaces newly added dirty manifests so
they are reviewed deliberately rather than slipping through. A
structural CI guard checks that the activation characterization
artifact regenerates without errors and contains all six expected
activation rows. The companion repository holds 122 unit and
integration tests that gate every push.

\section{Framework}
\label{sec:framework}

The infrastructure exists to answer the ``but what about\dots ?''
questions that ad-hoc CVNN comparisons in the literature routinely
sidestep.

\paragraph{CVNN library.}
We implement \code{ComplexLinear} (matmul + bias-add, deliberately
bypassing \code{F.linear} to sidestep its documented MPS gap),
\code{ComplexConv1d} (thin wrapper over \code{F.conv1d}, which natively
supports complex tensors on CPU/MPS/CUDA), and \code{ComplexDropout}
using a single shared real Bernoulli mask following Trabelsi
et~al.~\cite{trabelsi2018deep}. Six activations span the Liouville
trade-off: \code{CReLU}, \code{ZReLU}, \code{ModReLU}~\cite{arjovsky2016unitary},
\code{ComplexCardioid}, \code{Siglog}, and \code{ComplexTanh} as the
cautionary baseline. Initialization (Glorot~\cite{glorot2010understanding}
and Kaiming~\cite{he2015delving}) is complex-aware in both rectangular
and polar (Rayleigh-magnitude / uniform-phase) forms.

\paragraph{Activation characterization toolkit.}
For each candidate activation, we render a Cauchy--Riemann residual map
over the complex plane, compute summary statistics (median / p95 /
max), and measure two gradient norms: an MLP-scaffolded one and an
\emph{intrinsic} one ($|\partial L / \partial \bar{z}|$ for
$L = \sum |f(z)|^2$ on uniformly sampled $z$). The intrinsic measure
decouples the activation's contribution from any surrounding linear
layers --- without it, $\tanh$'s 855 MLP-gradient looks
$\sim$8$\times$ larger than its 108 intrinsic contribution.

\paragraph{Baseline matching.}
Every claim of ``complex beats real'' is backed by four families
running on the same task with the same shared search space:
\emph{naive real-stacked} (real net taking $(\mathrm{Re}, \mathrm{Im})$
as 2 channels; $\sim$2$\times$ the parameters of the complex layer);
\emph{matched-parameter real} (real net sized down so its real
\code{numel} equals the complex's, counting each complex slot as 2);
\emph{matched-FLOP real} (sized to match the complex MAdd count);
\emph{exact reparameterization} (the block-real $\bigl[\begin{smallmatrix}
A & -B \\ B & A \end{smallmatrix}\bigr]$ construction proven to
reproduce the complex forward bit-for-bit; a witness, not a baseline to
beat).

\paragraph{Budgeted random search.}
Every model family draws from the \emph{same} shared search space and
consumes the same trial $\times$ seed budget. The sweep harness samples
the hyperparameter list once and runs it against each family in turn.
The primary matched comparison chooses the complex reference trial by
mean validation accuracy, then reports every real baseline at that
same trial index. Independent family winners are also written as
diagnostics but are not the matched-capacity paper table.

\section{Findings}
\label{sec:findings}

\subsection{Activation characterization}
\label{sec:activation-characterization}

\begin{figure}[t]
  \centering
  \includegraphics[width=\linewidth]{figures/activation_tradeoff.png}
  \caption{Liouville trade-off, made empirical. \emph{Left}: gradient
    norm at init for each candidate activation, log scale.
    \emph{Right}: median and p95 Cauchy--Riemann residual on a
    $[-3, 3]^2$ grid in $\mathbb{C}$, log scale. \code{complex\_tanh}
    is locally holomorphic but unbounded (poles at $\pm i\pi/2$ inside
    the grid); \code{siglog} is bounded but never holomorphic;
    \code{crelu}, \code{zrelu}, \code{modrelu}, and
    \code{complex\_cardioid} sit between, each making a different
    compromise.}
  \label{fig:activation-tradeoff}
\end{figure}

\Cref{fig:activation-tradeoff} makes the Liouville bind empirically
visible. \code{complex\_tanh} has near-zero CR residual at the median
--- it really is locally holomorphic --- but the same metric's max is
905, because the grid (extent $= 3$) brackets the first pair of poles
at $z = \pm i\pi/2 \approx \pm 1.5708\,i$. The intrinsic Jacobian mean
is $\sim$108. By contrast, \code{siglog} is bounded ($|f| < 1$
everywhere) and never blows up, but its CR residual is everywhere
positive (p95 = 0.25) --- it is not holomorphic anywhere, by
construction. The remaining four activations sit between, each making
a different compromise.

This is not a finding about which activation is best in training ---
it is a visible measurement of the formal trade-off the project is
built around. The choice of activation in the downstream benchmarks
(default \code{crelu}) is consistent across runs and isolated from the
architectural variables.

\subsection{Synthetic phase classification --- null result}
\label{sec:phase-classification}

\begin{figure}[t]
  \centering
  \includegraphics[width=0.75\linewidth]{figures/synthetic_phase_classification_swept.png}
  \caption{1-D phase classification: matched shared-trial test
    accuracy after a 16$\times$3 sweep. All four families converge to
    $\sim$76\% with overlapping CIs --- the complex inductive bias has
    nothing to bite on at this scope.}
  \label{fig:phase-swept}
\end{figure}

After a 16$\times$3 budgeted sweep, the matched shared-trial comparison
puts all four families at $\sim$76.2--76.7\% test accuracy with
overlapping CIs (\Cref{fig:phase-swept}). On 1-D phase classification,
\textbf{the complex inductive bias does not help.} A 2-D real network
on $(\mathrm{Re}, \mathrm{Im})$ already has access to all the
information the task contains; there is no temporal or geometric
structure for $\mathbb{C}$-multiplication to exploit. The
\Cref{lem:phase-equivariance} of \Cref{sec:phase-equivariance} predicts
exactly this --- a single-sample phase task does not invoke any
sequence-level equivariance.

\subsection{Synthetic RF modulation --- complex wins under matched budget}
\label{sec:synthetic-rf}

\begin{table}[t]
  \centering
  \caption{Synthetic RF modulation classification, matched
    shared-trial selection after a 16$\times$6 budgeted sweep with the
    conv architecture. Complex beats every real baseline by
    $\geq$3.8~pp; CIs do not overlap.}
  \label{tab:synthetic-rf}
  \begin{tabular}{lrrrl}
    \toprule
    family & test acc & std & params & 95\% CI ($\sigma/\sqrt{n}$) \\
    \midrule
    \textbf{complex} & \textbf{0.8191} & 0.0163 & \textbf{3{,}974} & $[0.8061, 0.8322]$ \\
    real (stacked) & 0.7740 & 0.0226 & 2{,}099 & $[0.7559, 0.7920]$ \\
    real ($\approx$params) & 0.7769 & 0.0187 & 3{,}809 & $[0.7619, 0.7918]$ \\
    real ($\approx$FLOPs) & 0.7810 & 0.0195 & 7{,}779 & $[0.7654, 0.7966]$ \\
    \bottomrule
  \end{tabular}
\end{table}

\Cref{tab:synthetic-rf} reports the matched shared-trial comparison
after a 16$\times$6 sweep with the conv architecture
(\code{ComplexConv1d} $\to$ activation $\to$ \code{ComplexConv1d} $\to$
activation $\to$ global mean pool $\to$ \code{ComplexLinear} $\to$
$|\cdot|$). Complex beats every real baseline by $\geq$3.8~percentage
points; CIs do not overlap, and Welch two-sample $t$-tests give
$t \approx 4.2$, $\mathrm{df} \approx 9.8$ versus matched-parameter
real and $t \approx 3.7$, $\mathrm{df} \approx 9.7$ versus matched-FLOP
real ($p < 0.01$ in both cases).

\subsection{RadioML 2018.01A --- a positive result with a subtler story}
\label{sec:radioml}

The corrected\footnote{The first attempt at this sweep used the
synthetic benchmark's odd-stepped SNR list, which the loader silently
dropped because RadioML 2018.01A only ships even SNRs. The fix raised
on empty buckets by default; see the released code.}
\code{CReLU} run on a 3-class subset (BPSK / QPSK / 8PSK at
$\{-10, -6, -2, 2, 6, 10, 14, 18\}$~dB,
\code{max\_per\_class\_per\_snr=256}, \code{sample\_length=128}) on an
NVIDIA A100 produced a striking 23~pp gap:

\begin{table}[t]
  \centering
  \caption{RadioML 2018.01A subset, CReLU run, matched
    shared-trial selection after 16$\times$6 sweep. Read in isolation,
    a dramatic CVNN-vs-real win at matched parameters --- but see
    \Cref{sec:radioml-ablation} and \Cref{sec:radioml-mechanism}.}
  \label{tab:radioml-crelu}
  \begin{tabular}{lrrr}
    \toprule
    family & test acc & std & params \\
    \midrule
    \textbf{complex (CReLU)} & \textbf{0.7293} & 0.0085 & 58{,}886 \\
    real (stacked) & 0.4583 & 0.1394 & 29{,}891 \\
    real ($\approx$params) & 0.4999 & 0.1327 & 58{,}413 \\
    real ($\approx$FLOPs) & 0.4245 & 0.1415 & 117{,}123 \\
    \bottomrule
  \end{tabular}
\end{table}

Read in isolation, \Cref{tab:radioml-crelu} looks like a dramatic
complex-vs-real win at matched parameters --- and an earlier draft
claimed exactly that. The activation ablation in
\Cref{sec:radioml-ablation} forces a different reading.

\subsubsection{Activation ablation reveals a methodology-driven artifact}
\label{sec:radioml-ablation}

\begin{figure}[t]
  \centering
  \includegraphics[width=0.95\linewidth]{figures/radioml_activation_ablation.png}
  \caption{Five-activation ablation on the same RadioML scaffold.
    The CVNN line is nearly flat across all five activations
    (test accuracy $0.668$--$0.733$); real-baseline lines swing wildly
    between $\sim$0.45 and $\sim$0.70 depending on which activation
    the complex side uses. The 23~pp gap from
    \Cref{tab:radioml-crelu} is mostly real baselines collapsing under
    the matched-shared-trial selected hyperparameters, not complex
    pulling ahead.}
  \label{fig:radioml-ablation}
\end{figure}

We re-ran the same sweep four more times, varying only the complex
side's activation; the real baselines always use \code{ReLU}. The
matched-shared-trial selection per activation (\Cref{tab:ablation-summary})
shows complex's accuracy nearly flat (range 0.065) while
real baselines swing by 0.27. \Cref{fig:radioml-ablation} makes the
asymmetry visible. The 22--23~pp gap on \code{crelu} / \code{cardioid}
/ \code{siglog} \emph{is mostly the real baselines failing under
matched-shared-trial selection}, not the complex network pulling
ahead. The robustness asymmetry is the durable, defensible finding;
the apparently large per-activation gap is a consequence of that
asymmetry, not a separate phenomenon.

\begin{table}[t]
  \centering
  \caption{Five-activation ablation, matched shared-trial test
    accuracy. Under stable activations (\code{modrelu}, \code{zrelu})
    real baselines are also stable and the gap drops to $\pm 3$~pp;
    under unstable activations the gap inflates. See
    \Cref{sec:radioml-mechanism} for the mechanism.}
  \label{tab:ablation-summary}
  \begin{tabular}{lrrrrr}
    \toprule
    activation & complex & best real & gap (pp) & $\sigma$(complex) & mean $\sigma$(real) \\
    \midrule
    \code{crelu} & 0.7293 & 0.4999 & $+22.94$ & 0.009 & 0.138 \\
    \code{cardioid} & 0.7282 & 0.5028 & $+22.54$ & 0.019 & 0.134 \\
    \code{siglog} & 0.7014 & 0.4689 & $+23.25$ & 0.016 & 0.139 \\
    \code{modrelu} & 0.6683 & 0.6940 & $-2.58$ & 0.031 & 0.045 \\
    \code{zrelu} & 0.7330 & 0.7039 & $+2.91$ & 0.015 & 0.015 \\
    \bottomrule
  \end{tabular}
\end{table}

\subsubsection{Per-SNR breakdown}

At $-10$~dB everyone is near chance ($\sim$33\% for 3 classes); from
$+2$~dB onward, the \code{CReLU} complex network reaches 92.1\% at
$+10$~dB and 91.6\% at $+18$~dB while the real baselines plateau at
47--61\%. This is the same pattern as \Cref{tab:radioml-crelu}: the
gap reflects real baselines failing to train at the selected
hyperparameters under \code{CReLU}, not complex pulling cleanly away
under every stable activation regime.

\subsubsection{Comparison to published RadioML numbers}

The RadioML 2018.01A dataset paper~\cite{oshea2018radioml} reports a
real-valued ResNet-style classifier reaching roughly 90\% top-1 at
high SNR on the full 24-class task with the full sample length 1024.
Subsequent CVNN-on-RadioML
work~\cite{krzyston2020complex,tu2020complex} reports a typical
complex-vs-real advantage of a few percentage points on the same
24-class setup with comparable architectures. Our $+3$-and-flat
ablation result on \code{modrelu} / \code{zrelu} (the activations under
which real baselines are also stable) is consistent with that small
reported gap. The 22--23~pp gap on \code{crelu} / \code{cardioid} /
\code{siglog} is \emph{not} a contradiction of the literature --- it
is the matched-shared-trial selection rule magnifying real-baseline
instability that the literature's separate-tuning protocols would not
expose.

\subsubsection{Mechanism --- explosion-into-dead-region at step 1}
\label{sec:radioml-mechanism}

\Cref{sec:radioml-ablation} reports an asymmetry without explaining
it. We re-ran the matched-shared-trial selected configuration with
per-step instrumentation (train loss, total parameter gradient norm,
per-parameter gradient norm, max parameter magnitude) for each of
$5 \times 4 \times 3 = 60$ runs, capped at 200 steps (the divergence
pattern appears within the first 5).

\begin{figure}[t]
  \centering
  \includegraphics[width=0.95\linewidth]{figures/radioml_telemetry_grad.png}
  \caption{RadioML telemetry --- total gradient norm post-backward,
    log $y$-axis, one panel per activation. Real-baseline spikes at
    step 1 under \code{crelu} / \code{cardioid} / \code{siglog} mark
    the explosion-into-dead-region failure mode that produces the
    inflated gaps in \Cref{tab:ablation-summary}; under \code{modrelu}
    and \code{zrelu} no such spikes occur and the gap collapses to
    $\pm 3$~pp.}
  \label{fig:telemetry-grad}
\end{figure}

The mechanism is unambiguous (\Cref{fig:telemetry-grad}). On
\code{crelu} / \code{cardioid} / \code{siglog}, matched-shared-trial
picks a high-lr / wide-hidden config that complex selects because it
tolerates it (lr $\in \{0.024, 0.024, 0.040\}$, hidden = 64). At
step~1, every family sees a large AdamW step away from random init.
For real baselines this produces a step-1 train loss in $[3, 36]$, a
total gradient norm in $[6, 42]$, and a \code{head.weight} gradient in
$[6, 41]$ --- i.e.\ nearly all of the explosion sits on the final
classifier's weights. For the complex network the same step produces a
step-1 loss in $[1.07, 1.28]$, total gradient norm $[0.3, 1.0]$,
\code{head.weight} gradient $[0.1, 0.7]$ --- about 50--80$\times$
smaller than its real counterparts on identical data and learning rate.

The complex network's $(\mathrm{Re}, \mathrm{Im})$ parameter coupling
distributes the loss signal across more parameter slots, so the same
$\mathrm{lr} \times \mathrm{CE\text{-}gradient}$ product produces a
smaller effective step on \code{head.weight}. The real network does
not have that smoothing; AdamW's first step pushes some seeds into a
region where the second \code{Conv1d}'s ReLU activations are all-zero
(or all-active and saturated), gradients vanish, and training
``stabilizes'' at uniform-prediction chance level. We confirmed this:
\emph{the dead seeds end at exactly $\mathrm{loss} = \ln 3 \approx
1.099$} (3-class CE under uniform predictions), with
\code{total\_grad\_norm} $\leq 0.05$, and zero test-accuracy
improvement from random init.

\Cref{tab:dead-seeds} summarizes the dead-seed counts. \textbf{The
dead-seed rate jumps from 0\% to 33\% as soon as the
matched-shared-trial selected lr crosses $\sim$0.02.} Two orders of
magnitude separation in the step-1 head gradient between the unstable
and stable activation regimes.

\begin{table}[t]
  \centering
  \caption{Dead-seed counts (out of 9 per activation: three families
    $\times$ three seeds), step-1 \code{head.weight} gradient. The
    threshold behavior is striking: 0\% $\to$ 33\% as lr crosses
    $\sim$0.02. The same pattern replicates on synthetic AWGN data
    (\Cref{sec:synthetic-confirmation}).}
  \label{tab:dead-seeds}
  \begin{tabular}{lrrr}
    \toprule
    activation & lr & dead seeds & step-1 head.weight grad (real, max) \\
    \midrule
    \code{crelu} & 0.024 & $3/9$ & 13.7 \\
    \code{cardioid} & 0.024 & $3/9$ & 19.5 \\
    \code{siglog} & 0.040 & $3/9$ & 40.4 \\
    \code{modrelu} & 0.008 & $\mathbf{0/9}$ & 0.17 \\
    \code{zrelu} & 0.0024 & $\mathbf{0/9}$ & 0.13 \\
    \bottomrule
  \end{tabular}
\end{table}

\subsubsection{Cross-task confirmation}
\label{sec:synthetic-confirmation}

To check whether the explosion-into-dead-region pattern is
RadioML-specific (pulse shaping, carrier offset, channel effects) or
generic to the hp regime, we re-ran the same 60-run telemetry against
the synthetic AWGN-only RF generator using the same per-activation
hyperparameters that the RadioML matched-shared-trial selected.
\Cref{tab:cross-task} shows the side-by-side dead-seed counts.

\begin{table}[t]
  \centering
  \caption{Cross-task replication: same matched-shared-trial
    hyperparameters applied to synthetic AWGN-only RF data.
    Dead-seed counts replicate almost exactly; step-1 gradients on
    synthetic are if anything larger than RadioML's because AWGN-only
    signals carry more per-sample variance than pulse-shaped,
    channel-attenuated ones. \emph{The threshold is hp-driven, not
    data-driven.}}
  \label{tab:cross-task}
  \begin{tabular}{lrrrr}
    \toprule
    activation & RadioML dead/9 & synthetic dead/9 & RadioML max real grad & synthetic max real grad \\
    \midrule
    \code{crelu} & 3 & 3 & 19.9 & 34.1 \\
    \code{cardioid} & 3 & 3 & 19.9 & 34.1 \\
    \code{siglog} & 3 & 2 & 42.1 & 55.7 \\
    \code{modrelu} & $\mathbf{0}$ & $\mathbf{0}$ & 0.3 & 0.5 \\
    \code{zrelu} & $\mathbf{0}$ & $\mathbf{0}$ & 0.2 & 0.7 \\
    \bottomrule
  \end{tabular}
\end{table}

\section{Interpretation}
\label{sec:interpretation}

The three findings together support a more nuanced falsifiable story
than the original ``complex wins'' framing:

\begin{quote}
  The complex inductive bias pays for itself when the task carries
  structure that complex multiplication naturally encodes --- and is
  neutral otherwise. On real-data IQ classification, complex networks
  are \emph{robust to activation choice and seed} in a way
  capacity-matched real networks are not; this robustness asymmetry,
  rather than a raw accuracy advantage on a level playing field, is
  what produces large reported gaps under matched-shared-trial
  selection.
\end{quote}

\subsection{Why complex helps on IQ data --- phase equivariance}
\label{sec:phase-equivariance}

The intuition that ``complex multiplication has the right semantics
for phase-bearing signals'' can be made formal in one short lemma.

\begin{lemma}[Phase equivariance of \code{ComplexConv1d}]
  \label{lem:phase-equivariance}
  Let $C(z)[t] = \sum_{\tau} K[\tau] \cdot z[t-\tau] + b$ be a
  \code{ComplexConv1d} layer with complex-valued kernel $K$, complex
  bias $b$, and complex input sequence $z$. For any global phase
  $\varphi \in \mathbb{R}$ and any $t$,
  \[
    C(e^{i\varphi} \cdot z)[t]
    = e^{i\varphi} \cdot C(z)[t] + (1 - e^{i\varphi}) \cdot b.
  \]
\end{lemma}

\begin{proof}
  By linearity of complex multiplication and convolution,
  \[
    C(e^{i\varphi} \cdot z)[t]
    = \sum_{\tau} K[\tau] \cdot (e^{i\varphi} \cdot z[t-\tau]) + b
    = e^{i\varphi} \sum_{\tau} K[\tau] \cdot z[t-\tau] + b
    = e^{i\varphi} \cdot \bigl(C(z)[t] - b\bigr) + b. \qedhere
  \]
\end{proof}

In the bias-free case (or after subtracting the bias) the result is
exactly equivariant. The biased case adds a fixed offset
$(1 - e^{i\varphi}) \cdot b$ that the next layer can absorb; in our
networks the head's $|\cdot|$ magnitude operation is itself
phase-invariant ($|e^{i\varphi} \cdot y| = |y|$ for real $\varphi$), so
the \emph{output} logits are unaffected by a constant input phase
shift --- the architecture is phase-invariant as a whole at the
readout.

The same does not hold for \code{Conv1d} operating on a stacked
$(\mathrm{Re}\,z, \mathrm{Im}\,z)$ real input. A global phase shift
$e^{i\varphi}$ corresponds to multiplying each $(I, Q)$ sample by the
$2 \times 2$ rotation matrix $R_\varphi$. To match
\code{ComplexConv1d}'s equivariance, every real conv kernel would need
to commute with $R_\varphi$ applied to its 2 input channels --- a
constraint real \code{Conv1d} does not impose by construction. The
real network \emph{can} learn the symmetry from data, but it has to
spend capacity on it; the complex network gets it for free.

This is the structural inductive bias the synthetic phase
classification result (\Cref{sec:phase-classification}) and the
RadioML result (\Cref{sec:radioml}) probe in opposite regimes. A
single complex sample (phase classification) gives no temporal
structure for convolutional equivariance to bite on, and the 2-D real
network has the same ``see real and imag'' feature view, so the
complex network buys nothing --- the null is what
\Cref{lem:phase-equivariance} predicts. A length-$L$ IQ sequence
(synthetic RF, RadioML) does have temporal structure, the carrier
phase varies sample-to-sample under channel effects, and a
phase-equivariant conv stack matches that symmetry --- the wins are
also what the lemma predicts.

\paragraph{What the lemma does not explain.}
Phase equivariance of the \emph{layer} is necessary for the
architectural prior, but the activations between conv layers vary in
their own equivariance. By direct computation: \code{modrelu} and
\code{siglog} are phase-equivariant; \code{crelu}, \code{cardioid},
and \code{zrelu} are not. If the \Cref{sec:radioml} robustness
asymmetry were \emph{driven} by activation phase equivariance, the
stable activations should be exactly \code{modrelu} and \code{siglog}.
They are not: \Cref{sec:radioml-mechanism} shows \code{modrelu} and
\code{zrelu} are stable ($0/9$ dead seeds across nine real-baseline
runs each), while \code{crelu}, \code{cardioid}, and \code{siglog} are
unstable ($3/9$ dead). One phase-equivariant activation is on each
side of the divide. \textbf{Phase equivariance of the layer is
necessary for the inductive bias, but the robustness asymmetry
\Cref{sec:radioml-mechanism} measures lives one level lower --- at the
optimization dynamics induced by $(\mathrm{Re}, \mathrm{Im})$
parameter coupling on the classifier head, not at the equivariance of
the activation.} The two effects are partially independent and
partially compounding. This also explains why the
\Cref{sec:radioml} result is most defensible at activations like
\code{zrelu} where complex wins by a small margin rather than at
\code{crelu} where the apparent gap is amplified by the
optimization-side mechanism.

\section{Limitations}
\label{sec:limitations}

\begin{itemize}
  \item \textbf{RadioML subset, not the full archive.} The headline
    real-data run used BPSK / QPSK / 8PSK at 8 SNR levels with
    \code{max\_per\_class\_per\_snr=256}. The full archive has 24
    modulations $\times$ 26 SNRs and a much larger per-bucket count.
  \item \textbf{Three modulations only.} Angle-only constellations
    (BPSK, QPSK, 8PSK). RadioML \emph{does} have pulse-shaped QAM,
    so a separate run with larger conv stacks should separate them;
    not measured here.
  \item \textbf{Six seeds.} CIs are computed from a Gaussian
    approximation of the per-seed mean ($\sigma/\sqrt{n}$). A
    bootstrap CI on this many seeds gives similar widths but is not
    substantially more honest.
  \item \textbf{One loss.} Classification runs use cross-entropy.
  \item \textbf{No \code{ComplexBatchNorm} or \code{ComplexLayerNorm}.}
    Landing them would let us train deeper conv stacks fairly.
  \item \textbf{Modest hardware.} CPU and Apple Silicon for
    development; one A100 for the headline RF sweep. No
    multi-machine reproducibility check.
\end{itemize}

\section{Future work}
\label{sec:future-work}

\begin{itemize}
  \item \textbf{Scale up the RadioML run} to the full
    24-modulation $\times$ 26-SNR archive on the full 1024-sample
    length. The infrastructure is in place
    (\code{sweep\_radioml.py --preset full}); the run is deferred from
    the headline write-up to keep the paper anchored on the 3-class
    subset that supports the mechanism analysis. Recommended first
    activation is \code{zrelu} --- the only ablation point where real
    baselines also train cleanly, so the comparison is least confounded
    by the matched-shared-trial selection's interaction with the
    explosion mechanism.
  \item \code{ComplexBatchNorm} (whitening 2$\times$2 covariance, per
    Trabelsi et al.~\cite{trabelsi2018deep}) and deeper conv stacks.
  \item Audio (MUSDB18 STFT) and single-coil MRI reconstruction
    (fastMRI) --- the two natural 2-D scale-up targets.
  \item A separate study of per-activation training dynamics: does
    \code{ModReLU}'s bias-sign regime visibly affect convergence?
\end{itemize}

\section{Reproducibility}
\label{sec:reproducibility}

A C\&F submission's value is in what others can build on. Every claim
above is backed by a single CLI invocation against committed
artifacts:

\begin{verbatim}
# 1. Reproduce all three benchmarks
uv sync --all-groups
uv run python experiments/synthetic/sweep_phase_classification.py
uv run python experiments/rf/sweep_synthetic_modulation.py --device cuda
uv run python experiments/rf/sweep_radioml.py --device cuda \
    --activation crelu --seeds 0 1 2 3 4 5
# (repeat for cardioid / siglog / modrelu / zrelu)

# 2. Reproduce the mechanism analysis
uv run python scripts/run_gradient_telemetry.py
uv run python scripts/run_synthetic_rf_telemetry.py

# 3. Regenerate every figure from committed JSON
uv run python scripts/generate_paper_figures.py

# 4. Optional: scale up to the full RadioML archive
uv run python experiments/rf/sweep_radioml.py --device cuda \
    --preset full --activation zrelu --seeds 0 1 2
\end{verbatim}

Each result manifest records its \code{git\_commit}, environment, and
a \code{git\_dirty} flag computed on a tree that excludes
\code{results/} and the activation-characterization output directory
(so regen does not flip the bit on itself). CI surfaces newly added or
modified manifests with \code{git\_dirty: true} as warnings on the PR.
A structural CI guard checks the committed activation
characterization artifact contains all six expected activation rows
after a regen.

\bibliographystyle{plain}
\bibliography{references}

\end{document}
